\section{Exhibit A-2: protocol holonomy and route mismatch}
\label{sec:exhibit-holonomy}

This exhibit instantiates \Pthree (protocol / holonomy) as a quantitative diagnostic rather than a geometric claim.
The premise is simple: when there are multiple natural routes for constructing a packaged object from staged data, the routes need not commute at finite stage.
The discrepancy between routes is a defect; a stable closure should make that defect shrink under refinement.

\subsection{Two routes under coordinate change}
We consider a small coordinate change (a ``chart'' perturbation)
\[
\varphi(x) \;=\; x + \varepsilon x^2,
\qquad \varepsilon\ \text{small},
\]
and study a discrete derivative approximation at step size $h$.

There are (at least) two natural routes:
\begin{enumerate}
  \item \textbf{Route A (refine then differentiate then pull back).}
  Refine the sampled data, apply the discrete derivative in the target coordinate, then interpolate/pull back the result along $\varphi$.
  \item \textbf{Route B (pull back then refine then differentiate).}
  Pull back the sampled data along $\varphi$ at coarse scale, refine in the source coordinate, apply the discrete derivative, and convert by the chain rule factor $(\varphi')^{-1}$.
\end{enumerate}
At finite $h$, these routes generally disagree because interpolation, refinement, and discrete differentiation do not strictly commute.

\subsection{Route mismatch as a defect ledger}
In the \SBT dictionary, the existence of multiple admissible routes is \Pthree, while the numerical discrepancy between them is tracked by \Psix\ accounting.
We quantify the mismatch by a normalized route-mismatch statistic
\begin{equation}\label{eq:rm-def}
\mathrm{RM}(h)
\;:=\;
\frac{\|R_A(h)-R_B(h)\|}{\|R_A(h)\|+\varepsilon_0},
\end{equation}
where $R_A(h),R_B(h)$ are the route outputs on a shared comparison grid, $\|\cdot\|$ is an $\ell^2$ norm over evaluation points, and $\varepsilon_0>0$ is a small stabilizer.

A stable closure regime is indicated by \emph{decay} of $\mathrm{RM}(h)$ as $h\downarrow 0$.
Conversely, persistent or growing mismatch is evidence that the proposed packaging (or protocol family) is not yet feasible at the tested scales.

\subsection{Observed scaling under refinement}
In the recorded run, $\mathrm{RM}(h)$ decreases with refinement and is well fit by a power law
\[
\mathrm{RM}(h) \approx C\,h^{p}.
\]
The fitted exponent is
\[
p \;=\; \HolonomyExponent,
\]
where $\HolonomyExponent$ is imported from the repository artifacts (generated TeX macros).
Table~\ref{tab:holonomy} lists the measured values and Figure~\ref{fig:holonomy-rm} shows the decay trend.

As a null calibration, setting $\varepsilon=0$ (identity coordinate map) yields route mismatch at or below \HolonomyNullRMMax\ on the tested grids, confirming that the reported decay is not an interpolation bookkeeping artifact.
The log--log slope summary is reported together with the fit quality $R^2=\HolonomyFitRtwo$ and should be interpreted as a coarse scaling diagnostic rather than a precise convergence order.

% Auto-generated by scripts/export_results_tex.py from notes/*_last_run.json
\begin{table}[h]
\centering
\caption{Holonomy RM vs h}
\label{tab:holonomy}
\begin{tabular}{l|l}
\hline
 $h$ & RM \\
\hline
 \ensuremath{0.015625} & \ensuremath{0.00476218} \\
\hline
 \ensuremath{0.0078125} & \ensuremath{0.00132002} \\
\hline
 \ensuremath{0.00390625} & \ensuremath{4.9099\times 10^{-4}} \\
\hline
 \ensuremath{0.00195312} & \ensuremath{2.2498\times 10^{-4}} \\
\hline
\end{tabular}
\end{table}


\begin{figure}[htbp]
  \centering
  \includegraphics[width=0.85\linewidth]{holonomy_rm_last_run.png}
  \caption{Route mismatch under coordinate change: $\mathrm{RM}(h)$ versus step size $h$ for two noncommuting refinement routes (Equation~\ref{eq:rm-def}). The mismatch decays under refinement; the fitted exponent reported in the text is sourced from the generated macro $\HolonomyExponent$.}
  \label{fig:holonomy-rm}
\end{figure}

\subsection{Guardrail: holonomy is not directionality}
A crucial framework guardrail is that \emph{holonomy does not by itself certify directionality}.
In this instantiation, $\mathrm{RM}(h)$ is a measure of \emph{route dependence} (a \Pthree\ phenomenon) and its decay is a \emph{feasibility signal} (tracked by \Psix).
None of this implies an ``arrow of time'' or a monotone irreversible ordering on states.
It only says: given a staged protocol family and a proposed packaging rule, the protocol routes become increasingly compatible as resolution increases.

In later mathematical instantiations (e.g.\ differential geometry), true holonomy is tied to curvature and obstruction-to-patching phenomena.
Here we remain at the diagnostic level: route mismatch is a computable defect that should shrink in stable regimes, and can therefore be used to stress-test closure claims.
